\documentclass[11pt]{article}
\usepackage[T1]{fontenc}
\usepackage[utf8]{inputenc}
\usepackage{lmodern,amsmath,amssymb,microtype,xcolor}
\usepackage[margin=.9in]{geometry}
\providecommand{\tightlist}{\setlength{\itemsep}{0pt}\setlength{\parskip}{0pt}}
\usepackage[colorlinks=true,urlcolor=blue!60!black]{hyperref}
\setlength{\parindent}{0pt}
\setlength{\parskip}{.5em}
\setlength{\emergencystretch}{2em}
\author{Wentao Li\\\small Prepared with assistance from OpenAI Codex}
\date{November 22, 2026}
\title{Three Boolean summands do not imply the full Hindman property in ZF}
\hypersetup{pdfauthor={Wentao Li},pdftitle={Three Boolean summands do not imply the full Hindman property in ZF}}
\begin{document}
\maketitle
Brot and Cao described the missing implication in their 2018 Michigan REU report. Brot, Cao and Fernández-Bretón ask in Question 5.1(1) whether a model of ZFA can have an \(H\)-finite but \(H_3\)-infinite set. Their Proposition 4.15 constructs a model that works for colorings determined by cardinality but admits an obstructing structural coloring. The following permutation group gives a fresh component where \emph{every} supported coloring with a finite palette depends only on cardinality. A second kind of support makes finite petals of any required size available there. \href{https://lsa.umich.edu/content/dam/math-assets/math-document/reu-documents/2018FinalPapers/Brot\%20and\%20Cao.pdf}{Brot and Cao, 2018, p.~13} · \href{https://arxiv.org/abs/1910.11025}{Brot, Cao and Fernández-Bretón, 2019 preprint} · \href{https://doi.org/10.1016/j.apal.2021.102961}{2021 journal article}

\textbf{Theorem.} Relative to the consistency of ZFC, there is a model of pure ZF with an infinite set \(X\) such that \([X]^{<\omega}\) is Dedekind-finite, but every finite coloring of \([X]^{<\omega}\) admits an infinite family \(Y\) for which all symmetric differences of one, two or three distinct members of \(Y\) have one color. Theorem 3.2 of the source identifies Dedekind-finiteness of \([X]^{<\omega}\) with \(H\)-finiteness. In particular a binary coloring witnesses failure of the full Boolean Hindman property. The bound four for the known collapse to the full property is sharp.

\textbf{A locally finite highly transitive group.} Put \(D_j=\{0,\ldots,j+1\}\) and let \(\Omega\) consist of the eventually zero elements of \(\prod_{j<\omega}D_j\). The finite group \(G_m\) of all permutations of the length-\(m\) prefixes acts on \(\Omega\) by replacing the prefix and leaving the tail unchanged. Embed \(G_m\) diagonally in \(G_{m+1}\) and set \(G=\bigcup_m G_m\). A finite list of group elements generates a finite group. Any two ordered tuples of distinct points of equal length can be matched by an element of \(G\): take \(m\) large enough to distinguish their entries and beyond all nonzero digits of both tuples. For every \(k\ge1\) there is \(h\in G\) all of whose cycles on \(\Omega\) have size \(k\). Choose \(m\) for which \(k\) divides \(|D_0\times\cdots\times D_{m-1}|=(m+1)!\) and partition the prefixes into cycles of that length.

\textbf{Supports and finite orbits.} Use countably many disjoint, individually named copies of \(\Omega\) as atoms, with full product group acting componentwise. For the normal filter take finite intersections of point stabilizers and centralizers of single-component elements of \(G\), and all their supergroups. Conjugates of these generating subgroups are of the same form, and every atom has a support. This defines a Fraenkel--Mostowski model. In ordinary language, a support can name finitely many points and finitely many permutations, each acting on just one component. Naming \(h\) makes its graph and the partition into its cycles symmetric.

For a support let \(P\) be its finite set of points and let \(H\) be the subgroup generated by the named permutations. This is a finite subgroup acting on finitely many components. If \(z\notin H[P]\), it has infinite orbit under the subgroup fixing \(P\) pointwise and centralizing \(H\). To see this work in its component and pick \(m\) containing all named prefix permutations there, and beyond all nonzero digits of \(z\) and of the named points there. At a longer prefix length, the \(H\)-orbit of the prefix of \(z\) has arbitrarily many disjoint equivariantly isomorphic copies, obtained by varying the digits after \(m\). Swap the original orbit with one of these copies by that isomorphism and its inverse, fixing all other prefixes. This permutation commutes with the component projection of \(H\) and fixes every named point, since the original orbit misses \(P\) and the new orbit can be chosen to do so. The images of \(z\) give arbitrarily many distinct points. This proves the assertion, including in a component with no local parameters.

Let \(X\) be the set of all the point atoms. If \((F_n)_{n<\omega}\) is a supported sequence of finite subsets of \(X\), the stabilizer of a single support preserves each \(F_n\) setwise, since natural numbers are fixed. Every point of their union has finite orbit. Thus \(\bigcup_n F_n\subseteq H[P]\) is finite. In particular \([X]^{<\omega}\) has no countably infinite subset.

\textbf{The three-summand witness.} Fix a coloring \(c:[X]^{<\omega}\to q\) with \(q\) finite, and take a component untouched by its support. Arbitrary elements of \(G\) in this component fix the coloring and the pure color labels. Its restriction there has the form \(d(|s|)\) by high transitivity. Schur's theorem applied to \(j\mapsto d(2j)\) gives positive even \(u,v,u+v\) of one color. Interchange \(u,v\) to have \(u\ge v\), and put \(k=v/2\) and \(a=u-k\). Fix an \(a\)-point root \(R\) in the component and take one of its uniform \(k\)-cycle permutations \(h\). The family \[Y=\{R\cup O: O\text{ is an }h\text{-cycle disjoint from }R\}\] is infinite and supported by \(R\) and \(h\). The petals are pairwise disjoint. The sizes of the one-, two- and three-term symmetric differences are respectively \(a+k=u\), \(2k=v\) and \(a+3k=u+v\). The restriction of \(c\) therefore gives one color throughout. This Schur calculation is due to Brot, Cao and Fernández-Bretón. The new support structure makes their cardinality argument apply to arbitrary finite colorings.

\textbf{Transfer and scope.} The local conjunction transfers from ZFA to ZF by Jech--Sochor. Finite subsets of \(X\), their finite symmetric differences, families of finite subsets, finite-palette colorings and functions from \(\omega\) to \([X]^{<\omega}\) are all coded within a fixed finite number of power sets over \(X\cup\omega\) with tags. Finiteness and infinitude of the candidate families are tested by the same bounded codes. The argument needs no assumption on whether \(X\) itself can be linearly ordered. The proposed significance is \textbf{66/100}: it settles the missing implication in the authors' classification and locates the exact collapse at four summands. Independent proof and priority review are pending.
\end{document}
